% Owner: Kapish Rohilla. Proposed Solution: assembler and benchmark programs.
\subsection{Assembler and Benchmark Programs}
The Phase~I test programs are short (5 to 21 instructions) and were encoded by
hand, which once produced a hex file that did not match its assembly source. To
compare the pipeline with the baseline on realistic code, where CPI depends on
the dynamic instruction mix~\cite{kapishHennessyQA}, Phase~II adds a small
assembler and six benchmark programs with a larger instruction mix.

\subsubsection{Assembler}
\texttt{tools/assemble.py} is a two-pass assembler written in Python~3 with the
standard library only. The first pass removes \texttt{\#} comments, records the
byte address of every label and collects the instruction lines. The second pass
encodes each instruction into the R, I, S, B, J or U format of the RISC-V base
ISA~\cite{kapishRVSpec,kapishPattersonRV}. It accepts exactly the 29 instructions
of the processor, registers \texttt{x0}--\texttt{x31}, decimal and \texttt{0x}
immediates, labels, and the pseudo-instruction \texttt{nop}
(\texttt{addi x0, x0, 0}). A branch or \texttt{jal} target is a label or a
numeric byte offset relative to the instruction, as in the Phase~I sources.
Every other mnemonic, an unknown register, an undefined label or an immediate
that does not fit its field is rejected with the file name and line number. The
output has one word per line as eight hex digits followed by a
\texttt{//~source line} comment, which \texttt{\$readmemh} reads directly.
The option \texttt{--mix} counts the instructions of a program per class (ALU,
load, store, branch, jump, upper immediate).

The assembler was validated against the Phase~I programs: the option
\texttt{--selftest} assembles every \texttt{testbench/programs/test*\_*.s} file and
compares the result word by word with the committed \texttt{.hex} file. All six
programs (73 words in total) match.

\subsubsection{Benchmark Programs}
The six programs in Table~\ref{tab:bench} were written in assembly and
assembled with the tool. Each one builds its own input data (data memory starts
at zero), uses only the 29 supported instructions and data memory words 0 to
511, writes its results at byte address \texttt{0x100} and above, and ends with
the halt \texttt{jal x0, 0}. Together they cover the hazard cases of the
pipeline: dependent ALU chains (forwarding), a loaded value used by the next
instruction (load-use stall), and taken branches and jumps (flush).

\begin{table*}[!t]
\centering
\caption{Benchmark Programs}
\label{tab:bench}
\setlength{\tabcolsep}{4pt}
\renewcommand{\arraystretch}{1.1}
\begin{tabular}{@{}lp{5.2cm}p{4.0cm}p{5.4cm}r@{}}
\toprule
\textbf{Program} & \textbf{Computes} & \textbf{Why chosen} & \textbf{Hazards exercised} & \textbf{$N$} \\ \midrule
\texttt{bench\_fib}    & Iterative Fibonacci $F(n)$ for $n = 1, 10, 20, 30, 40, 46$ in a subroutine &
  Tight ALU loop, function calls & Back-to-back ALU dependences (forwarding), taken loop branch, JAL/JALR flushes & 785 \\
\texttt{bench\_sum}    & Fills a 32-element array with $13i+5$ and sums it (6608) &
  Simple array traversal & Load-use stall on every element, taken loop branch & 330 \\
\texttt{bench\_sort}   & Bubble sort of 16 signed words generated by xorshift32 &
  Mix of loads, stores and data-dependent branches & Load-use stall before each compare, branches taken or not taken depending on the data, store after load & 920 \\
\texttt{bench\_gcd}    & Subtraction-based GCD of six pairs, e.g.\ $\gcd(1764,1260)=252$ &
  Branch-heavy control flow & Two branches per iteration with operands forwarded into the branch unit, JAL/JALR flushes & 786 \\
\texttt{bench\_memcpy} & Copies 64 words, then the sum and XOR of the copy &
  Load/store-heavy kernel & Store data taken from the preceding load and checksum add after a load (load-use stalls) & 1231 \\
\texttt{bench\_mul}    & Shift-and-add product (low 32 bits) of six pairs, including negative operands &
  The subset has no MUL & ANDI result used by the next BEQ, data-dependent branches, LUI + ADDI constants & 570 \\
\bottomrule
\end{tabular}
\end{table*}

$N$ in Table~\ref{tab:bench} is the number of instructions executed up to and
including the halt, measured by the benchmark testbench; it is the same on both
designs. Table~\ref{tab:mix} gives the static instruction mix of each program
as printed by \texttt{python3 tools/assemble.py --mix}. The jump column includes
the final halt.

\begin{table}[!t]
\centering
\caption{Static Instruction Mix of the Benchmarks (\texttt{--mix})}
\label{tab:mix}
\setlength{\tabcolsep}{3.5pt}
\renewcommand{\arraystretch}{1.08}
\begin{tabular}{@{}lccccccc@{}}
\toprule
\textbf{Program} & \textbf{ALU} & \textbf{Load} & \textbf{Store} & \textbf{Branch} & \textbf{Jump} & \textbf{LUI} & \textbf{Total} \\ \midrule
fib    & 14 & 0 & 6 & 2 & 8  & 0 & 30 \\
sum    & 13 & 1 & 3 & 2 & 1  & 0 & 20 \\
sort   & 16 & 2 & 3 & 4 & 1  & 1 & 27 \\
gcd    & 15 & 0 & 6 & 2 & 10 & 0 & 33 \\
memcpy & 23 & 2 & 4 & 3 & 1  & 1 & 34 \\
mul    & 19 & 0 & 6 & 2 & 9  & 5 & 41 \\ \bottomrule
\end{tabular}
\end{table}

\subsubsection{Self-Checking Benchmark Testbench}
One testbench, \texttt{testbench/tb\_bench.v}, runs every benchmark on both
designs: with \texttt{-DSINGLE} it instantiates \texttt{riscv\_single\_cycle\_top},
otherwise \texttt{riscv\_pipeline\_top}; the program and its name are given with
\texttt{-DPROG} and \texttt{-DNAME}. The expected result words were computed
independently with short Python scripts that implement each algorithm directly
(for example \texttt{math.gcd} and Python's \texttt{sorted}), never from the
processor's output, and are compared with \texttt{dut.u\_dmem.mem} after the run.
For \texttt{bench\_memcpy} all 64 destination words are also compared with the
source words. The simulation runs for a fixed 30\,000 cycles and records the
cycle in which the halt is reached; if the halt is never reached the run fails.
The pipeline counters use the same definitions as the Phase~I testbench
(instructions leaving ID, load-use stall cycles, taken branches and jumps, halt
cycle plus three), so the numbers are comparable with Phase~I. Each run prints
one line,
\texttt{BENCH name=\dots{} design=\dots{} instr=\dots{} cycles=\dots{} stalls=\dots{} taken=\dots{} result=PASS|FAIL},
and \texttt{simulation/run\_bench.sh} runs the twelve combinations and exits
with an error if any of them fails. All twelve runs pass, and every pipeline
cycle count satisfies $C = N + 4 + S + 2T$.
